\documentclass[11pt]{article}
\usepackage{amsmath,amssymb,amsthm}
\usepackage{algorithm}
\usepackage{algpseudocode}
\usepackage{geometry}
\geometry{margin=1in}
\newtheorem{theorem}{Theorem}
\newtheorem{lemma}{Lemma}
\newcommand{\E}{\mathbb{E}}
\newcommand{\R}{\mathbb{R}}
\newcommand{\norm}[1]{\left\lVert #1 \right\rVert}
\newcommand{\inner}[1]{\left\langle #1 \right\rangle}
\begin{document}

\section*{Algorithm Name}
\textbf{Noisy Gradient Descent (DP‑SGD with Gaussian Noise).}

\section*{Mathematical Formulation}
Let $f:\R^{d}\rightarrow\R$ be the objective.  
At iteration $t=0,1,\dots,T-1$ we compute a (possibly clipped) stochastic gradient  
\[
g_{t}= \nabla f(x_{t}) + \xi_{t},\qquad 
\xi_{t}\sim\mathcal N\!\bigl(0,\sigma^{2} I_{d}\bigr),
\]
and update
\begin{equation}
\boxed{\,x_{t+1}=x_{t}-\eta_{t}\,g_{t}\,} \label{eq:update}
\end{equation}
where $\eta_{t}>0$ is the stepsize.  
The Gaussian noise $\xi_{t}$ provides $(\varepsilon,\delta)$‑DP under the standard
Gaussian mechanism; the cumulative privacy loss is tracked via the moments accountant
or Rényi DP (not shown here).

\section*{Pseudocode}
\begin{algorithm}[H]
\caption{Noisy Gradient Descent (DP‑SGD)}\label{alg:dp-sgd}
\begin{algorithmic}[1]
\Require Learning‑rate schedule $\{\eta_{t}\}_{t=0}^{T-1}$, noise scale $\sigma$, clipping norm $C$ (optional)
\State Initialize $x_{0}\in\R^{d}$
\For{$t=0$ \textbf{to} $T-1$}
    \State Compute true gradient $g^{\text{true}}_{t}= \nabla f(x_{t})$
    \State (Optional) clip: $g^{\text{clip}}_{t}= g^{\text{true}}_{t}\cdot\min\!\Bigl(1,\frac{C}{\norm{g^{\text{true}}_{t}}}\Bigr)$
    \State Sample noise $\xi_{t}\sim\mathcal N(0,\sigma^{2}I_{d})$
    \State Form noisy gradient $g_{t}= g^{\text{clip}}_{t}+ \xi_{t}$
    \State Update $x_{t+1}=x_{t}-\eta_{t} g_{t}$   \Comment{Eq.~\eqref{eq:update}}
\EndFor
\State \textbf{return} $x_{T}$
\end{algorithmic}
\end{algorithm}

\section*{Dry Run Example}
Consider the one‑dimensional quadratic
\[
f(w)=(w-3)^{2},\qquad \nabla f(w)=2(w-3).
\]
Set $w_{0}=0$, constant stepsize $\eta=0.1$, and use \emph{zero} noise for illustration
($\xi_{t}=0$).  

\begin{center}
\begin{tabular}{c|c|c|c}
$t$ & $\nabla f(w_{t})$ & $w_{t+1}=w_{t}-\eta\,\nabla f(w_{t})$ & $f(w_{t+1})$\\ \hline
0 & $2(0-3)=-6$ & $0-0.1(-6)=0.6$ & $(0.6-3)^{2}=5.76$\\
1 & $2(0.6-3)=-4.8$ & $0.6-0.1(-4.8)=1.08$ & $(1.08-3)^{2}=3.6864$\\
2 & $2(1.08-3)=-3.84$ & $1.08-0.1(-3.84)=1.464$ & $(1.464-3)^{2}=2.3660$\\
\end{tabular}
\end{center}
The iterates move toward the optimum $w^{*}=3$, and the objective value decreases.

\newpage
\section*{Assumptions}
\begin{enumerate}
    \item \textbf{Convexity.} $f$ is convex:
    \[
    f(y)\ge f(x)+\inner{\nabla f(x)}{y-x},\qquad\forall x,y\in\R^{d}.
    \]
    \item \textbf{Smoothness.} $f$ is $L$‑smooth:
    \[
    \norm{\nabla f(x)-\nabla f(y)}\le L\norm{x-y},\qquad\forall x,y\in\R^{d}.
    \]
    Equivalently,
    \[
    f(y)\le f(x)+\inner{\nabla f(x)}{y-x}+\frac{L}{2}\norm{y-x}^{2}.
    \]
    \item \textbf{Bounded (clipped) gradients.} After optional clipping,
    \[
    \norm{\nabla f(x)}\le G\qquad\forall x.
    \]
    \item \textbf{Noise properties.} The added Gaussian noise satisfies
    \[
    \E[\xi_{t}\mid x_{t}]=0,\qquad 
    \E\!\bigl[\norm{\xi_{t}}^{2}\mid x_{t}\bigr]=d\sigma^{2},
    \]
    and is independent across iterations.
    \item \textbf{Stepsize.} The schedule $\{\eta_{t}\}$ is non‑increasing and satisfies
    $0<\eta_{t}\le \frac{1}{L}$ for all $t$.
\end{enumerate}

\section*{Main Convergence Theorem}
\begin{theorem}[Expected Suboptimality of Noisy GD]\label{thm:conv}
Under Assumptions 1–5, let $x^{*}$ be any minimizer of $f$.
Define the averaged iterate $\bar x_{T}= \frac{1}{T}\sum_{t=0}^{T-1}x_{t}$.
If a constant stepsize $\eta_{t}\equiv\eta\le\frac{1}{L}$ is used, then
\[
\boxed{
\E\!\bigl[f(\bar x_{T})-f(x^{*})\bigr]\;
\le\;
\frac{\norm{x_{0}-x^{*}}^{2}}{2\eta T}
\;+\;
\frac{\eta}{2}\bigl(G^{2}+d\sigma^{2}\bigr).
}
\]
Choosing $\eta = \frac{\norm{x_{0}-x^{*}}}{\sqrt{T\,(G^{2}+d\sigma^{2})}}$ (which respects
$\eta\le 1/L$ for sufficiently large $T$) yields the rate
\[
\E\!\bigl[f(\bar x_{T})-f(x^{*})\bigr]=
O\!\Bigl(\frac{1}{\sqrt{T}}\Bigr)
\;+\;O\!\Bigl(\frac{\sigma^{2}}{\sqrt{T}}\Bigr).
\]
\end{theorem}

\section*{Theoretical Properties}
\begin{itemize}
    \item \textbf{Stability.} The algorithm remains stable as long as $\eta\le 1/L$,
    because the smoothness inequality guarantees a descent step in expectation.
    \item \textbf{Hyper‑parameter sensitivity.}
    The bound shows a linear dependence on $\eta$ for the variance term
    $(G^{2}+d\sigma^{2})$, while the bias term decays as $1/(\eta T)$.
    Hence the optimal $\eta$ balances these two effects, giving the $1/\sqrt{T}$ rate.
    \item \textbf{Comparison to vanilla SGD.}
    Setting $\sigma=0$ recovers the classical SGD bound
    $\frac{\norm{x_{0}-x^{*}}^{2}}{2\eta T}+ \frac{\eta G^{2}}{2}$.
    The additional $d\sigma^{2}$ term quantifies the privacy‑induced noise penalty.
    \item \textbf{Privacy–utility trade‑off.}
    Larger $\sigma$ yields stronger $(\varepsilon,\delta)$‑DP guarantees but
    inflates the variance term, degrading the convergence rate.
\end{itemize}

\section*{Discussion}
The theorem demonstrates that DP‑SGD with Gaussian noise retains the
$O(1/\sqrt{T})$ convergence of standard SGD up to an additive term proportional
to the noise variance.  Because the bound holds for the \emph{averaged}
iterate, one can also obtain a high‑probability guarantee via Markov’s
inequality and standard concentration for sub‑Gaussian noise.
The result is tight for convex smooth objectives and matches known lower
bounds for differentially private optimisation.

\newpage
\section*{Preliminaries (Definitions and Lemmas)}
\begin{lemma}[Smoothness inequality]\label{lem:smooth}
If $f$ is $L$‑smooth, then for any $x,y\in\R^{d}$,
\[
f(y)\le f(x)+\inner{\nabla f(x)}{y-x}+\frac{L}{2}\norm{y-x}^{2}.
\]
\end{lemma}
\begin{lemma}[Non‑expansiveness of projection (optional)]\label{lem:proj}
If a clipping operator $\Pi_{C}$ satisfies $\norm{\Pi_{C}(v)}\le C$, then
$\norm{\Pi_{C}(v)-\Pi_{C}(u)}\le \norm{v-u}$.
\end{lemma}
Both lemmas are standard; proofs are omitted for brevity.

\section*{Main Proof (Step‑by‑Step Derivation)}
We prove Theorem~\ref{thm:conv}.  
All expectations are taken w.r.t. the randomness of the noise $\{\xi_{t}\}$.

\noindent\textbf{Step 1.} Expand the squared distance after one update.
\begin{align}
\norm{x_{t+1}-x^{*}}^{2}
&= \norm{x_{t}-\eta g_{t}-x^{*}}^{2} \nonumber\\
&= \norm{x_{t}-x^{*}}^{2}
-2\eta\inner{g_{t}}{x_{t}-x^{*}}
+\eta^{2}\norm{g_{t}}^{2}. \tag{Step 1}
\end{align}

\noindent\textbf{Step 2.} Substitute $g_{t}= \nabla f(x_{t})+\xi_{t}$ and take conditional expectation.
Using $\E[\xi_{t}\mid x_{t}]=0$,
\begin{align}
\E\!\bigl[\norm{x_{t+1}-x^{*}}^{2}\mid x_{t}\bigr]
&= \norm{x_{t}-x^{*}}^{2}
-2\eta\inner{\nabla f(x_{t})}{x_{t}-x^{*}}
+\eta^{2}\E\!\bigl[\norm{\nabla f(x_{t})+\xi_{t}}^{2}\mid x_{t}\bigr] \nonumber\\
&= \norm{x_{t}-x^{*}}^{2}
-2\eta\inner{\nabla f(x_{t})}{x_{t}-x^{*}}
+\eta^{2}\Bigl(\norm{\nabla f(x_{t})}^{2}
+\E\!\bigl[\norm{\xi_{t}}^{2}\mid x_{t}\bigr]\Bigr) \nonumber\\
&\le 
\norm{x_{t}-x^{*}}^{2}
-2\eta\inner{\nabla f(x_{t})}{x_{t}-x^{*}}
+\eta^{2}\bigl(G^{2}+d\sigma^{2}\bigr). \tag{Step 2}
\end{align}
The inequality uses the gradient bound $\norm{\nabla f(x_{t})}\le G$ and the noise variance
$\E[\norm{\xi_{t}}^{2}]=d\sigma^{2}$.

\noindent\textbf{Step 3.} Apply convexity to replace the inner product.
Convexity gives
\[
f(x_{t})-f(x^{*})\le \inner{\nabla f(x_{t})}{x_{t}-x^{*}}.
\]
Hence,
\[
-2\eta\inner{\nabla f(x_{t})}{x_{t}-x^{*}}
\le -2\eta\bigl(f(x_{t})-f(x^{*})\bigr). \tag{Step 3}
\]

\noindent\textbf{Step 4.} Combine Steps 2 and 3.
\begin{align}
\E\!\bigl[\norm{x_{t+1}-x^{*}}^{2}\mid x_{t}\bigr]
&\le 
\norm{x_{t}-x^{*}}^{2}
-2\eta\bigl(f(x_{t})-f(x^{*})\bigr)
+\eta^{2}\bigl(G^{2}+d\sigma^{2}\bigr). \tag{Step 4}
\end{align}

\noindent\textbf{Step 5.} Take total expectation and rearrange.
\begin{align}
2\eta\,\E\!\bigl[f(x_{t})-f(x^{*})\bigr]
&\le 
\E\!\bigl[\norm{x_{t}-x^{*}}^{2}\bigr]
-\E\!\bigl[\norm{x_{t+1}-x^{*}}^{2}\bigr]
+\eta^{2}\bigl(G^{2}+d\sigma^{2}\bigr). \tag{Step 5}
\end{align}

\noindent\textbf{Step 6.} Sum (\ref{Step 5}) over $t=0,\dots,T-1$.
Telescoping the distance terms yields
\begin{align}
2\eta\sum_{t=0}^{T-1}\E\!\bigl[f(x_{t})-f(x^{*})\bigr]
&\le 
\norm{x_{0}-x^{*}}^{2}
-\E\!\bigl[\norm{x_{T}-x^{*}}^{2}\bigr]
+T\eta^{2}\bigl(G^{2}+d\sigma^{2}\bigr) \nonumber\\
&\le 
\norm{x_{0}-x^{*}}^{2}
+T\eta^{2}\bigl(G^{2}+d\sigma^{2}\bigr). \tag{Step 6}
\end{align}
The non‑negativity of the norm term gives the last inequality.

\noindent\textbf{Step 7.} Divide by $2\eta T$ and use convexity of $f$ to replace the average.
Because $f$ is convex,
\[
f(\bar x_{T})\le \frac{1}{T}\sum_{t=0}^{T-1}f(x_{t}).
\]
Thus,
\begin{align}
\E\!\bigl[f(\bar x_{T})-f(x^{*})\bigr]
&\le 
\frac{\norm{x_{0}-x^{*}}^{2}}{2\eta T}
+\frac{\eta}{2}\bigl(G^{2}+d\sigma^{2}\bigr). \tag{Step 7}
\end{align}
This is exactly the bound stated in Theorem~\ref{thm:conv}.

\noindent\textbf{Step 8.} Optimising $\eta$.  
Set the derivative of the right‑hand side of (\ref{Step 7}) w.r.t. $\eta$ to zero:
\[
-\frac{\norm{x_{0}-x^{*}}^{2}}{2\eta^{2}T}
+\frac{1}{2}\bigl(G^{2}+d\sigma^{2}\bigr)=0
\;\Longrightarrow\;
\eta^{\star}= \frac{\norm{x_{0}-x^{*}}}{\sqrt{T\,(G^{2}+d\sigma^{2})}}.
\]
Plugging $\eta^{\star}$ into (\ref{Step 7}) gives
\[
\E\!\bigl[f(\bar x_{T})-f(x^{*})\bigr]
\le
\frac{\norm{x_{0}-x^{*}}\,\sqrt{G^{2}+d\sigma^{2}}}{\sqrt{T}}
=O\!\Bigl(\frac{1}{\sqrt{T}}\Bigr)
+O\!\Bigl(\frac{\sigma^{2}}{\sqrt{T}}\Bigr),
\]
which completes the proof. \hfill $\square$

\section*{Rate Derivation (With Constants)}
From Step 7 we have the explicit constant‑level bound
\[
\E\!\bigl[f(\bar x_{T})-f(x^{*})\bigr]
\le
\underbrace{\frac{\norm{x_{0}-x^{*}}^{2}}{2\eta T}}_{\text{bias term}}
\;+\;
\underbrace{\frac{\eta}{2}\bigl(G^{2}+d\sigma^{2}\bigr)}_{\text{variance term}}.
\]
Choosing $\eta=1/L$ (the largest stepsize allowed by smoothness) yields
\[
\E\!\bigl[f(\bar x_{T})-f(x^{*})\bigr]
\le
\frac{L\norm{x_{0}-x^{*}}^{2}}{2T}
+\frac{G^{2}+d\sigma^{2}}{2L}.
\]
If $T$ is large, the first term dominates and we obtain the classic $O(1/T)$
behaviour for strongly convex objectives (not covered here).  For the convex case,
the $1/\sqrt{T}$ rate derived with the optimal $\eta$ is tight.

\section*{Special Cases}
\begin{itemize}
    \item \textbf{No privacy noise ($\sigma=0$).} The bound reduces to the standard
    SGD result $\frac{\norm{x_{0}-x^{*}}^{2}}{2\eta T}+\frac{\eta G^{2}}{2}$.
    \item \textbf{Zero gradient bound ($G=0$).} When the true gradient vanishes
    (e.g., at the optimum), the error is purely due to noise:
    $\frac{\norm{x_{0}-x^{*}}^{2}}{2\eta T}+\frac{\eta d\sigma^{2}}{2}$,
    minimised at $\eta^{\star}\propto 1/\sqrt{T}$ giving $O(\sigma/\sqrt{T})$.
    \item \textbf{High‑dimensional regime.} The variance term scales linearly with
    dimension $d$, reflecting the fact that Gaussian noise is added to each
    coordinate.  Dimensionality reduction or per‑coordinate clipping can mitigate
    this effect.
\end{itemize}

\end{document}